\documentclass[a4paper,11pt]{article}

\usepackage[utf8]{inputenc}
\usepackage[T1]{fontenc}
\usepackage{newtxtext}
\usepackage[left=2.5cm, right=2.5cm, top=2.5cm, bottom=2.5cm]{geometry}
\usepackage[french]{babel}
\usepackage{setspace}
\usepackage{graphicx}
\usepackage{etoolbox}
\usepackage{csquotes}
\usepackage{caption}
\usepackage{fancyhdr}
\usepackage{titlesec}
\usepackage{tocloft}
\usepackage{amsmath,amssymb,amsfonts}
\usepackage{mathrsfs}
\usepackage{float}
\usepackage{dsfont}
\usepackage{tikz}
\usepackage{url}
\usepackage{color, xcolor}
\usepackage{array}
\usepackage{booktabs}
\usepackage{colortbl}
\usepackage{longtable}
\usepackage{pdflscape}
\usepackage{enumitem}
\usepackage{mdframed}
\usepackage{hyperref}
\usetikzlibrary{arrows.meta, positioning, calc, shapes.geometric, fit, backgrounds}

% ---------------------------------------------------------------- charte ----
% Charte reprise d'outil_USP (docs/latex/CONVENTIONS.md, section 2) : bleu
% financier pour la structure, vert pour l'accent et les sorties, deux gris de
% tableau.
\definecolor{maincolor}{RGB}{0,70,140}
\definecolor{greenline}{RGB}{0,180,80}
\definecolor{tableheader}{RGB}{200,220,250}
\definecolor{tablerow}{RGB}{245,245,245}
\definecolor{financeblue}{RGB}{0,70,140}

\newcommand{\annee}{2026}
\newcommand{\datedoc}{\today}
\newcommand{\versiondoc}{0.1}
\newcommand{\statutdoc}{Projet}
% Titre court utilise dans le pied de page (fancyhdr).
\newcommand{\titredoc}{Documentation de l'outil de construction de tables d'expérience}

\captionsetup[figure]{labelfont={bf, color=maincolor}, textfont={normalfont}, labelsep=colon}
\captionsetup[table]{labelfont={bf, color=maincolor}, textfont={normalfont}, labelsep=colon}

\hypersetup{colorlinks=true, linkcolor=financeblue, urlcolor=financeblue,
            pdftitle={Documentation de l'outil de construction de tables d'experience}}

\renewcommand{\baselinestretch}{1.08}
\pagestyle{fancy}
\fancyhf{}
\fancyhead[L]{\color{maincolor}\nouppercase{\leftmark}}
\renewcommand{\headrulewidth}{0pt}
\renewcommand{\footrulewidth}{0.4pt}
\fancyfoot[R]{\color{greenline}\textbf{page \thepage}}
\fancyfoot[L]{\color{maincolor}\titredoc}
\fancypagestyle{plain}{\fancyhf{}%
  \renewcommand{\headrulewidth}{0pt}\renewcommand{\footrulewidth}{0.4pt}%
  \fancyfoot[R]{\color{greenline}\textbf{page \thepage}}%
  \fancyfoot[L]{\color{maincolor}\titredoc}}

% ------------------------------------------------------ titres de section ---
% Le bandeau bleu pleine largeur avec le titre en vert marque le niveau
% section, qui joue ici le role du chapitre.
\newcommand{\titrebandeau}[1]{%
  \noindent{\color{financeblue}\rule{\linewidth}{2.5pt}}
  \par\vspace{5mm}
  {\color{greenline}\raggedright\fontsize{19}{23}\bfseries\selectfont #1\par}
  \vspace{6mm}}

\titleformat{\section}[display]
  {\normalfont}{}{0pt}
  {\titrebandeau}
\titlespacing*{\section}{0pt}{0pt}{0pt}
\titleformat{\subsection}
  {\color{financeblue}\normalfont\large\bfseries}{\thesubsection}{0.7em}{}
\titleformat{\subsubsection}
  {\color{financeblue}\normalfont\normalsize\bfseries}{\thesubsubsection}{0.6em}{}

\renewcommand{\cftsecfont}{\bfseries\color{financeblue}}
\renewcommand{\cftsecaftersnum}{\color{financeblue}}
\renewcommand{\cftsecpagefont}{\color{financeblue}}
\renewcommand{\cftsubsecfont}{\color{black}}
\renewcommand{\cftsubsecpagefont}{\color{black}}
\setlength{\cftsecnumwidth}{2.2em}
\setcounter{tocdepth}{2}

\newcommand{\code}[1]{\texttt{\small #1}}
\newcommand{\refl}[1]{{\small\textcolor{maincolor!70!black}{\textit{Référence.} #1}}\par\smallskip}

% Entree d'un catalogue lu par la batterie de concordance
% (tests/concordance_doc_moteur.R ; docs/latex/CONVENTIONS.md, section 3.2) :
% #1 = nom du catalogue (exports, fonctions-<fichier>, roles, recettes...),
% non imprime ; #2 = nom tel qu'il s'ecrit dans le code, compose comme \code.
\newcommand{\catalogue}[2]{\code{#2}}

% Encadre de citation d'un texte qui fait foi : filet vertical bleu en marge.
\newcommand{\reglement}[2]{%
  \par\vspace{2mm}\noindent\begin{minipage}{\linewidth}
  \small\setlength{\parindent}{0pt}
  \textcolor{financeblue}{\vrule width 2pt}\hspace{4mm}%
  \begin{minipage}{\dimexpr\linewidth-7mm\relax}
  \textcolor{financeblue}{\textbf{#1}}
  \par\vspace{2pt}
  {\itshape #2}
  \end{minipage}
  \end{minipage}\par\vspace{3mm}}

% L'encadre generique reprend le filet bleu de la charte.
\newmdenv[linecolor=financeblue, linewidth=0pt, leftline=true,
          innerleftmargin=6mm, innertopmargin=4pt, innerbottommargin=4pt,
          backgroundcolor=tablerow, skipabove=3mm, skipbelow=3mm]{encadre}

% Titre de partie : un seul filet bleu, titre en vert dessous, corps plus
% petit que le bandeau des sections qui suivent (charte d'outil_USP).
\newcommand{\partie}[2]{%
  \clearpage
  \vspace*{1cm}
  \noindent{\color{financeblue}\rule{\linewidth}{2pt}}\par\vspace{3mm}
  {\color{financeblue}\raggedright\normalsize\bfseries\scshape #1\par}
  \vspace{1.5mm}
  {\color{greenline}\raggedright\fontsize{20}{24}\bfseries\selectfont #2\par}
  \vspace{5mm}
  \addcontentsline{toc}{section}{\protect\numberline{}\textsc{#1} --- #2}}

\emergencystretch=1.5em

\begin{document}

% ============================================================ page de garde ==
\thispagestyle{empty}
\begin{tikzpicture}[remember picture, overlay]
  \fill[financeblue]
    ([xshift=2.5cm,yshift=-3.2cm]current page.north west) rectangle
    ([xshift=-2.5cm,yshift=-3.28cm]current page.north east);
  \fill[greenline]
    ([xshift=2.5cm,yshift=-3.2cm]current page.north west) rectangle
    ([xshift=7.5cm,yshift=-3.28cm]current page.north west);
  \fill[financeblue]
    ([xshift=2.5cm,yshift=3.6cm]current page.south west) rectangle
    ([xshift=-2.5cm,yshift=3.53cm]current page.south east);
\end{tikzpicture}%
\vspace*{2.1cm}
\noindent
{\color{financeblue}\normalsize\bfseries\scshape Exercice \annee}

\vspace{1.7cm}
\noindent
{\color{financeblue}\fontsize{28}{34}\bfseries\selectfont
 Documentation\\ de l'outil\par}

\vspace{1.1cm}
\noindent
{\color{black}\Large
 Construction de tables d'expérience\\
 Package \code{tablexp}, méthodes, vérifications et application\par}

\vspace{0.7cm}
\noindent
{\color{black}\normalsize\itshape
 Une table d'expérience est l'estimation d'une transition d'un modèle multi-états\par}

\vfill
\noindent
{\footnotesize\color{black}
 Établie le \datedoc\quad\textbar\quad Version \versiondoc\quad\textbar\quad
 Statut : \statutdoc\quad\textbar\quad Package \code{tablexp} 0.0.0.9000\par}
\vspace{2.6cm}
\clearpage

% ================================================================= resume ==
\titrebandeau{Objet du document}
\markboth{Objet du document}{}

\noindent
Ce document est la documentation de l'\textbf{outil de construction de tables
d'expérience} \code{outil\_tables}. Il couvre le package R \code{tablexp}, qui
porte toute la logique quantitative, les vérifications qui l'accompagnent
(batteries de tests, contrôles du dépôt, plateforme de calcul) et, quand elle
existera, l'application interactive qui l'exploite. Une \emph{table
d'expérience} est construite à partir des observations d'un portefeuille
d'assurés~: pour un risque donné, l'outil estime le taux d'une
\emph{transition} (vivant $\to$ décédé, en cours $\to$ racheté\ldots) selon un
ou plusieurs axes, le lisse ou le positionne sur une table de référence, le
ferme aux âges extrêmes, le valide, et restitue la table avec tout ce qui
permet de la rejouer et de la relire.

\medskip\noindent
Les destinataires des tables et des rapports sont, par priorité, les
hypothèses biométriques de la meilleure estimation (\emph{best estimate})
Solvabilité~II, que l'ACPR revoit en contrôle, puis la certification d'une
table d'expérience par un actuaire indépendant agréé (\code{docs/exigences.md},
E1.2). Le document doit permettre à un relecteur externe, sans lire le code,
de savoir ce qui a été calculé, avec quelles hypothèses, et de rejouer le
calcul (E3.7). Les textes qui font foi sont recensés par le cahier des charges
(\code{docs/exigences.md}, §~2.4)~; ils ne sont cités dans le présent document
qu'une fois établis sur la source (issue~\#13).

\medskip\noindent
Il est organisé en une partie par module du package (E3.1). La
\textbf{partie~I} présente l'outil~: forme du package et graphe des modules,
services transverses, reproductibilité, batteries de vérification et
plateforme de calcul, contrôles du dépôt. La \textbf{partie~II} décrit le point
d'entrée unique et le plan de construction. Les parties suivantes suivent
l'ordre des modules du moteur~: données, exposition, taux bruts, lissage,
positionnement, table prospective, fermeture, validation, rapport, puis le
simulateur, qui produit les jeux de données sous une loi vraie connue. La
dernière partie rassemble l'index des fonctions exportées, la traçabilité des
exigences et les références. Chaque partie de module se termine par le
tableau des fonctions du module, qui relie la description de la méthode à
son implémentation~; un module qui n'existe pas encore a déjà sa partie, avec
un tableau vide.

\medskip\noindent
Le document ne porte aucun chiffre issu d'une exécution (E3.4)~: les
résultats d'une construction sont dans son rapport figé et sa fiche de
construction, et les bilans des vérifications dans la sortie des batteries.
Ses conventions de rédaction sont fixées par \code{docs/latex/CONVENTIONS.md}~;
le vocabulaire est celui du glossaire du projet (\code{CONTEXT.md}).

\vspace{6mm}
\noindent
\begin{center}
\small
\rowcolors{2}{tablerow}{white}
\begin{tabular}{@{}p{4.6cm}p{10.2cm}@{}}
\rowcolor{tableheader}
\textcolor{financeblue}{\textbf{Élément}} & \textcolor{financeblue}{\textbf{Référence}} \\
Cahier des charges & \code{docs/exigences.md} (exigences E1.1 à E5.5) \\
Décisions d'architecture & \code{docs/adr/} (ADR 0002 à 0008) \\
Moteur de calcul & package \code{tablexp} (\code{R/}), packages de priorité
\code{base} seulement (ADR 0002) \\
Interface & \code{app/app.R}, \code{app/affichage.R} (à venir~; aucun calcul) \\
Version de R & 4.3.3~; \code{DESCRIPTION} déclare \code{Depends: R (>= 4.3.3)} \\
Plateforme de référence & conteneur \code{rocker/r-ver:4.3.3} épinglé par
empreinte, BLAS et LAPACK de référence (ADR 0008~; section~\ref{sec:plateforme}) \\
Données & aucune donnée réelle~: jeux simulés sous une loi vraie et une graine \\
\end{tabular}
\end{center}

\clearpage
\titrebandeau{Sommaire}
\markboth{Sommaire}{}
% \tableofcontents composerait son propre titre, qui ferait doublon avec le
% bandeau ci-dessus : on n'appelle que \@starttoc.
\makeatletter
\@starttoc{toc}
\makeatother
\newpage


% =============================================================================
\partie{Partie I}{Présentation de l'outil}
% =============================================================================

% =============================================================================
\section{Architecture du package et graphe des modules}
\label{sec:architecture}
% =============================================================================

\subsection{Principe de séparation}
\label{sec:separation}

L'outil sépare strictement le \emph{calcul} de l'\emph{affichage}. La règle
est la suivante~: si un résultat peut être calculé indépendamment de
l'interface, son calcul figure dans le package \code{tablexp} (E4.3). Le moteur
reçoit un \code{data.frame} et la déclaration des \emph{rôles} de ses colonnes,
jamais un fichier~; la lecture des fichiers, et les conversions de type qui
l'accompagnent, relèvent d'une couche mince hors du moteur, qui n'existe pas
encore (E4.4). Le moteur refuse un type incompatible, il ne le convertit pas.

\begin{center}
\small
\rowcolors{2}{tablerow}{white}
\begin{tabular}{@{}p{3.6cm}p{5.6cm}p{5.6cm}@{}}
\rowcolor{tableheader}
\textcolor{financeblue}{\textbf{Emplacement}} &
\textcolor{financeblue}{\textbf{Contenu}} &
\textcolor{financeblue}{\textbf{Interdit}} \\
\code{R/} (package \code{tablexp}) & toute la logique quantitative~: rôles et
validation des données, recettes d'exposition, estimation, lissage,
positionnement, prospective, fermeture, validation, simulation, quantités des
rapports et empreintes & toute primitive Shiny (\code{input\$},
\code{reactive}, \code{render*})~; tout package hors priorité \code{base}
sans ADR de méthode~; toute dépendance à Java \\
\code{app/app.R} (à venir) & déclaration des rôles, édition du plan, appel du
point d'entrée sur clic du bouton «~Relancer les calculs~», affichage &
tout calcul~; tout appel d'une fonction non exportée (\code{tablexp:::}) \\
\code{app/affichage.R} (à venir) & formatage, tableaux HTML, tracés à partir
des quantités du résultat & tout calcul \\
\code{tests/} & batteries, harnais, scripts de contrôle & --- \\
\end{tabular}
\end{center}

\medskip
L'application vit hors du package, sous \code{app/}, et non dans un fichier
\code{R/affichage.R}~: tout fichier de \code{R/} est installé dans l'espace de
noms du package (section~\ref{sec:fichiers-package}), si bien qu'un fichier
d'affichage placé là entrerait dans le moteur et rendrait circulaire le
contrôle de frontière (ADR 0007, point 1.3). Le dossier \code{app/} et son
script de frontière n'existent pas encore~: ils viennent avec la branche de
l'application.

\subsection{Fichiers du package}
\label{sec:fichiers-package}

Le package \code{tablexp} est à la racine du dépôt (ADR 0007, point 1.1)~:

\begin{itemize}[leftmargin=*, itemsep=2pt]
  \item \code{DESCRIPTION}~: \code{Package: tablexp}, \code{Depends: R (>= 4.3.3)}
    (décision M44), aucun champ \code{Imports}, \code{Suggests} ni
    \code{LinkingTo}, \code{Encoding: UTF-8}, \code{License: file LICENSE}~; le
    contrôle de la liste blanche (section~\ref{sec:controles-depot}) en fait une
    règle vérifiée~;
  \item \code{NAMESPACE}, écrit à la main~: une directive \code{export} par
    fonction, sans \code{exportPattern}~; les fonctions exportées sont les
    points d'entrée et les fonctions que l'application appelle (contrôles de
    saisie à sens métier, constructeurs d'objets partagés), recensées dans
    l'index (section~\ref{sec:index-exports})~; aucune méthode S3 n'est
    exportée~;
  \item \code{R/}~: un fichier par module et deux fichiers transversaux
    (tableau~\ref{tab:fichiers-r})~;
  \item \code{.Rbuildignore}~: exclut de l'archive du package \code{app/},
    \code{docs/}, \code{tests/}, \code{.claude/}, \code{.github/} et les
    fichiers de pilotage du dépôt.
\end{itemize}

\noindent
Aucun fichier \code{Rd} ni \code{roxygen2}~: la documentation de référence est
le présent document, tenu concordant avec le code par la batterie de
concordance (section~\ref{sec:batteries}~; E3.3).

\medskip
\textbf{\code{R/} ne contient que du code du package.} Une installation depuis
l'arbre de travail (\code{R CMD INSTALL <dossier>}), qui est celle des
batteries, n'applique pas \code{.Rbuildignore}~: tout fichier de \code{R/} est
installé, même s'il y est exclu (constat~: un fichier \code{R/b.R} exclu par
\code{.Rbuildignore} est présent dans le package installé~; ADR 0007, mesures).
Un fichier étranger au moteur ne peut donc pas être rangé dans \code{R/}.

\begin{center}
\small
\rowcolors{2}{tablerow}{white}
\begin{longtable}{@{}p{4.2cm}p{7.6cm}p{2.9cm}@{}}
\caption{Fichiers de \code{R/}, dans l'ordre des modules (ADR 0007, points 1.2 et 2.1).}
\label{tab:fichiers-r}\\
\rowcolor{tableheader}
\textcolor{financeblue}{\textbf{Fichier}} &
\textcolor{financeblue}{\textbf{Contenu}} &
\textcolor{financeblue}{\textbf{État}} \\
\endfirsthead
\rowcolor{tableheader}
\textcolor{financeblue}{\textbf{Fichier}} &
\textcolor{financeblue}{\textbf{Contenu}} &
\textcolor{financeblue}{\textbf{État}} \\
\endhead
\code{R/outils.R} & services transverses sans contenu métier~: graine locale,
refus et avertissements (section~\ref{sec:outils})~; n'appelle aucun module &
présent \\
\code{R/donnees.R} & rôles, contrôles, refus et exclusions, unité
d'observation, table de référence (section~\ref{sec:donnees}) & présent \\
\code{R/exposition.R} & recettes d'exposition, conventions d'âge, troncature
et censure (section~\ref{sec:exposition}) & présent \\
\code{R/brut.R} & taux bruts et intervalles (section~\ref{sec:brut}) & à créer \\
\code{R/lissage.R} & lissages (section~\ref{sec:lissage}) & à créer \\
\code{R/positionnement.R} & positionnement sur une table de référence
(section~\ref{sec:positionnement}) & à créer \\
\code{R/prospective.R} & tables prospectives (section~\ref{sec:prospective}) & à créer \\
\code{R/fermeture.R} & fermeture aux âges extrêmes (section~\ref{sec:fermeture}) & à créer \\
\code{R/validation.R} & validation et suivi d'expérience (section~\ref{sec:validation}) & à créer \\
\code{R/rapport.R} & quantités des rapports, fiche de construction, empreintes
(section~\ref{sec:rapport}) & présent \\
\code{R/simulation.R} & simulateur (section~\ref{sec:simulation})~; hors de la
chaîne du moteur & présent \\
\code{R/construire\_table.R} & point d'entrée unique, validation du plan,
orchestration seule, aucun calcul (section~\ref{sec:construction}) & présent \\
\end{longtable}
\end{center}

\subsection{Graphe d'appels autorisés}
\label{sec:graphe}

Les modules sont ordonnés~: \code{outils}, \code{donnees}, \code{exposition},
\code{brut}, \code{lissage}, \code{positionnement}, \code{prospective},
\code{fermeture}, \code{validation}, \code{rapport}. Un module ne peut appeler
que des fonctions définies dans les modules situés \textbf{strictement avant
lui} dans cet ordre (ADR 0007, point 2.1). Deux fichiers sont hors de la
chaîne~: \code{simulation}, qui ne peut appeler que \code{donnees} et
\code{outils} et qu'\textbf{aucun module du moteur n'appelle}, et
\code{construire\_table}, qui peut appeler tous les modules \textbf{sauf}
\code{simulation}. La figure~\ref{fig:graphe-modules} et le
tableau~\ref{tab:appels-autorises} donnent ce graphe.

\medskip
Deux raisons fondent cette règle. Un ordre sans cycle rend chaque module
lisible et vérifiable avec ses seuls modules amont~: une anomalie d'un module
ne peut atteindre que les modules situés après lui. Et le simulateur, qui
produit les jeux sur lesquels le moteur est testé, ne partage pas le code
qu'il sert à tester~: s'il pouvait appeler \code{exposition}, une erreur de la
recette d'exposition se retrouverait à l'identique dans les données simulées
et dans leur estimation, et la vérification ne la verrait pas (ADR 0007,
«~Options écartées~»). Le simulateur reproduit le même modèle, pas le même
code. Une seule exception à la lecture «~en aval seulement~» est assumée~:
\code{positionnement} peut appeler \code{lissage}, parce que la méthode~4 de
positionnement est un lissage de la table du moment avec dérive de la
référence (E2.12)~; le sens inverse est interdit, \code{lissage} n'appelant
jamais \code{positionnement}.

\begin{figure}[H]
\centering
\begin{tikzpicture}[
  font=\small,
  ent/.style   = {rectangle, rounded corners=2pt, draw=financeblue, line width=0.7pt,
                  fill=tableheader, text=financeblue, align=center, inner sep=3pt,
                  minimum height=6mm, text width=36mm, font=\small\bfseries},
  fonc/.style  = {rectangle, rounded corners=2pt, draw=financeblue, line width=0.7pt,
                  fill=white, text=black, align=center, inner sep=3pt,
                  minimum height=6mm, text width=36mm},
  prevu/.style = {rectangle, rounded corners=2pt, draw=financeblue!60, dashed,
                  line width=0.6pt, fill=white, text=black!60, align=center,
                  inner sep=3pt, minimum height=6mm, text width=36mm},
  fl/.style    = {-{Latex[length=2mm]}, draw=financeblue, line width=0.6pt},
  flx/.style   = {font=\scriptsize\itshape, text=greenline, inner sep=1.5pt}
]
\node[fonc] (rap) {\code{rapport}};
\node[prevu, below=3.5mm of rap] (val) {\code{validation}};
\node[prevu, below=3.5mm of val] (fer) {\code{fermeture}};
\node[prevu, below=3.5mm of fer] (pro) {\code{prospective}};
\node[prevu, below=3.5mm of pro] (pos) {\code{positionnement}};
\node[prevu, below=3.5mm of pos] (lis) {\code{lissage}};
\node[prevu, below=3.5mm of lis] (bru) {\code{brut}};
\node[fonc, below=3.5mm of bru] (exp) {\code{exposition}};
\node[fonc, below=3.5mm of exp] (don) {\code{donnees}};
\node[fonc, below=3.5mm of don] (out) {\code{outils}};

\node[ent, above=8mm of rap] (ct) {\code{construire\_table}\\[1pt]
  \scriptsize\mdseries point d'entrée, orchestration};
\node[fonc, right=22mm of don] (sim) {\code{simulation}\\[1pt]
  \scriptsize simulateur};

\foreach \a/\b in {rap/val, val/fer, fer/pro, pro/pos, pos/lis, lis/bru,
                   bru/exp, exp/don, don/out}
  \draw[fl] (\a) -- (\b);
\draw[fl] (ct) -- node[flx, right] {tous les modules du moteur} (rap);
\draw[fl] (sim) -- (don);
\end{tikzpicture}
\caption{Graphe d'appels autorisés des modules (ADR 0007, point 2.1). Une
flèche de $A$ vers $B$ signifie que $A$ peut appeler $B$ et, la règle étant
«~strictement avant~», tout module situé après $B$ dans le sens des flèches.
Blocs pleins~: fichiers présents dans \code{R/}~; blocs en pointillé~: modules
prévus, fichiers absents. \code{simulation} n'est appelé par aucun module du
moteur, ni par \code{construire\_table}~; \code{lissage} n'appelle jamais
\code{positionnement}.}
\label{fig:graphe-modules}
\end{figure}

\begin{center}
\small
\rowcolors{2}{tablerow}{white}
\begin{longtable}{@{}p{3.6cm}p{11.2cm}@{}}
\caption{Appels autorisés, par module appelant (ADR 0007, point 2.1~;
constante \code{APPELS\_AUTORISES} de \code{tests/controles/graphe\_modules.R}).}
\label{tab:appels-autorises}\\
\rowcolor{tableheader}
\textcolor{financeblue}{\textbf{Module appelant}} &
\textcolor{financeblue}{\textbf{Peut appeler}} \\
\endfirsthead
\rowcolor{tableheader}
\textcolor{financeblue}{\textbf{Module appelant}} &
\textcolor{financeblue}{\textbf{Peut appeler}} \\
\endhead
\code{outils} & aucun module (fonctions des packages de priorité \code{base} seulement) \\
\code{donnees} & \code{outils} \\
\code{exposition} & \code{donnees}, \code{outils} \\
\code{brut} & \code{exposition} et les modules amont \\
\code{lissage} & \code{brut} et les modules amont \\
\code{positionnement} & \code{lissage} et les modules amont \\
\code{prospective} & \code{positionnement} et les modules amont \\
\code{fermeture} & \code{prospective} et les modules amont \\
\code{validation} & \code{fermeture} et les modules amont \\
\code{rapport} & tous les modules ci-dessus \\
\code{simulation} & \code{donnees} et \code{outils} seulement \\
\code{construire\_table} & tous les modules sauf \code{simulation} \\
\end{longtable}
\end{center}

\subsubsection*{Contrôle par script}

Le graphe est vérifié par \code{tests/controles/graphe\_modules.R}, en R de
base (\code{codetools}, de priorité \code{recommended}, est exclu par
l'ADR 0002), dans le job «~Tests (R 4.3.3)~» de la CI (ADR 0007, point 2.2)~:

\begin{enumerate}[leftmargin=*, itemsep=2pt]
  \item chaque fichier \code{R/<module>.R} est lu par \code{parse()}~; un fichier
    dont le nom n'est pas un module connu est un écart~;
  \item les affectations de premier niveau (\code{nom <- \ldots} ou
    \code{nom = \ldots}, fonctions comme constantes) donnent la table
    «~symbole $\to$ module~»~; un symbole défini dans deux modules est un
    écart~;
  \item pour chaque définition, les noms qu'elle emploie (\code{all.names()}),
    croisés avec cette table, donnent les arêtes entre modules~; toute arête
    absente du tableau~\ref{tab:appels-autorises} est un écart, imprimé avec le
    module appelant, le module appelé et les deux symboles~;
  \item les appels dynamiques interdits (ci-dessous) sont cherchés dans tout
    l'arbre syntaxique.
\end{enumerate}

\noindent
Le script imprime toutes les arêtes trouvées et tous les écarts, puis sort
avec le code~1 s'il en a trouvé au moins un. Ses limites sont écrites dans son en-tête~: un nom local homonyme d'un
symbole d'un autre module compte comme une arête (faux positif possible,
jamais d'arête manquée par homonymie)~; un code de premier niveau qui n'est
pas une affectation est un écart~; les appels vers les packages de base ne
sont pas contrôlés ici, mais par la liste blanche de \code{DESCRIPTION}. Le
script est lui-même testé sur des témoins (\code{tests/unitaires/test\_controles.R})~:
il échoue, entre autres, quand \code{exposition} appelle \code{brut}, quand
\code{donnees} appelle \code{simulation}, quand \code{lissage} appelle
\code{positionnement} et quand \code{construire\_table} appelle
\code{simulation}, et il passe sur un graphe conforme où \code{positionnement}
appelle \code{lissage}.

\subsubsection*{Appels dynamiques interdits}

Un appel dont la cible est calculée à l'exécution (\code{do.call("f", \ldots)},
\code{get("f")}) n'apparaît pas dans \code{all.names()}~: il créerait une
arête invisible pour le contrôle. Le script interdit donc dans \code{R/}
(ADR 0007, point 2.3, liste élargie à la mise en œuvre et consignée dans
l'annotation du 16/10/2026)~:

\begin{itemize}[leftmargin=*, itemsep=2pt]
  \item \code{do.call()} dont l'argument \code{what}, apparié par
    \code{match.call()} (position ou nom), n'est ni un symbole ni une fonction
    anonyme (chaîne, \code{paste0(\ldots)}\ldots)~; \code{do.call(f, args)} avec
    \code{f} un symbole désignant une fonction reste permis~;
  \item toute occurrence, appel ou fonction passée comme objet
    (\code{lapply(x, get)}), de \code{match.fun}, \code{get}, \code{get0},
    \code{mget}, \code{getFromNamespace}, \code{getExportedValue},
    \code{getFunction}, \code{asNamespace} et \code{getNamespace}~;
  \item \code{eval()} ou \code{evalq()} d'un \code{parse()}, \code{str2lang()},
    \code{str2expression()}, \code{call()} ou \code{as.call()}~;
  \item l'opérateur \code{:::} et toute référence qualifiée \code{tablexp::}~;
  \item les formes qualifiées \code{base::}, \code{utils::} et \code{methods::}
    de ce qui précède, lues comme la forme nue.
\end{itemize}

\noindent
Une fonction passée en argument puis appelée dans le corps de la fonction qui la reçoit crée une arête
que le script attribue à l'appelant qui l'a passée~: c'est voulu, la
dépendance existe bien là. Les tests, seuls, appellent les fonctions non
exportées par \code{tablexp:::} (ADR 0007, point 11.3).

% =============================================================================
\section{Services transverses~: \code{R/outils.R}}
\label{sec:outils}
% =============================================================================

\code{R/outils.R} porte les services sans contenu métier partagés par tous les
modules~: l'exécution sous graine locale (section~\ref{sec:reproductibilite}),
la forme des refus et celle des avertissements (ADR 0007, point 1.2). Il
n'appelle aucun module. Sans lui, la fonction de graine locale serait dupliquée
dans le simulateur et dans le futur bootstrap par têtes, et chaque module
inventerait sa forme de refus.

\subsection{Refus~: une condition R typée}
\label{sec:refus}

Le moteur ne devine pas~: quand une entrée ne permet pas de construire, il
\emph{refuse} (E2.3, E2.7). Un refus est une \textbf{condition R typée, jamais
un objet retourné} (ADR 0007, point 6.4). \code{refus\_tablexp()} construit la
condition, de classe \code{c("refus\_tablexp", "error", "condition")}, avec
quatre champs~:

\begin{itemize}[leftmargin=*, itemsep=2pt]
  \item \code{module}~: le module qui refuse~;
  \item \code{motif}~: un code court stable, en minuscules ASCII, chiffres et
    soulignés, conforme au motif \code{MOTIF\_CODE\_REFUS}
    (\verb|^[a-z][a-z0-9_]*$|)~;
  \item \code{message}~: la phrase qui dit ce qui manque (E5.5)~;
  \item \code{details}~: une liste (rôle, colonne, étape, paramètre, décompte,
    premiers numéros de ligne).
\end{itemize}

\noindent
\code{refuser()} la lève par \code{stop()} et ne rend jamais la main. Un
argument mal formé de \code{refus\_tablexp()} (module ou message vide, motif
non conforme, détails qui ne sont pas une liste) n'est pas un refus mais une
erreur de programmation du package, signalée par \code{stop()}. Les tests
vérifient le motif, le décompte et les numéros par \code{tryCatch()}~;
l'application attrapera la condition et affichera son message. Une condition
est préférée à un objet portant \code{ok = FALSE}, forme d'\code{outil\_USP}~:
un tel objet peut être lu par mégarde comme une construction partielle, une
condition ne le peut pas. Les codes de refus de chaque module sont recensés
dans la partie du module.

\subsection{Avertissements métier}
\label{sec:avertissements}

\emph{À rédiger (issue~\#9).}

\subsection{Fonctions du module}
\label{sec:fonctions-outils}

{\small
\rowcolors{2}{tablerow}{white}
\begin{longtable}{@{}p{4.2cm}p{1.6cm}p{3.4cm}p{5.4cm}@{}}
\caption{Définitions de premier niveau de \code{R/outils.R}.}
\label{tab:fonctions-outils}\\
\rowcolor{tableheader}
\textcolor{financeblue}{\textbf{Nom}} &
\textcolor{financeblue}{\textbf{Statut}} &
\textcolor{financeblue}{\textbf{Entrées}} &
\textcolor{financeblue}{\textbf{Sortie et rôle}} \\
\endfirsthead
\rowcolor{tableheader}
\textcolor{financeblue}{\textbf{Nom}} &
\textcolor{financeblue}{\textbf{Statut}} &
\textcolor{financeblue}{\textbf{Entrées}} &
\textcolor{financeblue}{\textbf{Sortie et rôle}} \\
\endhead
\catalogue{fonctions-outils}{GENERATEUR} & constante & --- & générateur fixé
(section~\ref{sec:generateur})~: Mersenne-Twister, inversion pour la loi
normale, rejet pour \code{sample()} \\
\catalogue{fonctions-outils}{MOTIF\_CODE\_REFUS} & constante & --- & motif
d'un code de refus ou d'avertissement~: \verb|^[a-z][a-z0-9_]*$| \\
\catalogue{fonctions-outils}{refus\_tablexp()} & interne & \code{module},
\code{motif}, \code{message}, \code{details} (liste, vide par défaut) &
condition de classe \code{refus\_tablexp}, sans la lever~; \code{stop()} si un
argument est mal formé (section~\ref{sec:refus}) \\
\catalogue{fonctions-outils}{refuser()} & interne & mêmes arguments & lève le
refus~; ne rend jamais la main \\
\catalogue{fonctions-outils}{avertissements\_vides()} & interne & --- &
accumulateur vide des avertissements métier
(section~\ref{sec:avertissements}~; issue~\#9) \\
\catalogue{fonctions-outils}{avertir()} & interne & accumulateur, \code{etape},
\code{code}, \code{message}, \code{nombre} & accumulateur augmenté d'une ligne,
sans \code{warning()} (section~\ref{sec:avertissements}~; issue~\#9) \\
\catalogue{fonctions-outils}{sous\_graine()} & interne & \code{graine} (entier
explicite, sans défaut), \code{expr} & valeur de \code{expr} évaluée sous la
graine et \code{GENERATEUR}, état du générateur de l'appelant restauré~; refus
\code{graine\_invalide} (section~\ref{sec:sous-graine}) \\
\end{longtable}
}

% =============================================================================
\section{Reproductibilité et aléa}
\label{sec:reproductibilite}
% =============================================================================

\subsection{Exigence et portée}
\label{sec:reproductibilite-exigence}

À données, rôles, plan et graine identiques, le programme produit des
résultats identiques \textbf{au bit près sur une même machine}~; toute
simulation passe par une graine locale explicite sous un générateur fixé,
consignée dans le résultat, l'état global du générateur étant restauré à la
sortie~; aucun horodatage n'entre dans la partie du résultat comparée lors
d'un rejeu (E4.6). La présente section décrit les deux mécanismes du socle~:
le générateur fixé et l'exécution sous graine locale. La consignation de la
graine et du générateur dans la fiche de construction et le rejeu de la fiche
sont décrits avec la fiche (section~\ref{sec:fiche}).

\begin{encadre}
\small
\textbf{Statut.} L'identité au bit près est un \emph{constat} des batteries
(section~\ref{sec:verifications-alea}), établi sur la machine qui les exécute~;
ce n'est pas une propriété démontrée. L'identité entre deux machines n'est
\textbf{pas garantie}~: deux processeurs différents pourraient différer au
dernier bit, même dans le même conteneur, si la bibliothèque mathématique
du système choisit ses variantes de calcul selon le processeur (hypothèse de
l'ADR 0008, non mesurée). C'est pourquoi les références de non-régression
seront comparées à tolérance et produites sur une plateforme unique
(section~\ref{sec:plateforme}).
\end{encadre}

\subsection{Générateur fixé}
\label{sec:generateur}

Le générateur est la constante \code{GENERATEUR} de \code{R/outils.R}~:
\code{kind = "Mersenne-Twister"}, \code{normal.kind = "Inversion"},
\code{sample.kind = "Rejection"}. Ce sont les valeurs par défaut de R 4.3.3
(\code{RNGkind()} dans une session neuve), fixées en constante pour qu'un
changement de défaut d'une version de R ne change pas les tirages en silence.
Le générateur uniforme est le Mersenne-Twister de Matsumoto et Nishimura~; les
lois normales s'obtiennent par inversion de la fonction de répartition~; la
génération uniforme discrète (\code{sample()}) emploie la méthode par rejet,
défaut de R depuis la version 3.6.0, la méthode antérieure par arrondi rendant
\code{sample()} sensiblement non uniforme sur de grandes populations
(documentation de R, \code{?Random}).

\refl{M.~Matsumoto, T.~Nishimura, \emph{Mersenne Twister: a 623-dimensionally
equidistributed uniform pseudo-random number generator}, \emph{ACM
Transactions on Modeling and Computer Simulation}, 8, 1998, p.~3--30 (cité par
la documentation de R, \code{?Random}).}

\noindent
Un changement de \code{GENERATEUR} change tout résultat simulé~: c'est un
changement de résultat, soumis au visa (ADR 0007, «~Conséquences~»).

\subsection{Exécution sous graine locale}
\label{sec:sous-graine}

\code{sous\_graine(graine, expr)} est la seule voie par laquelle le package
tire un aléa (ADR 0007, point 8.1). Elle procède ainsi~:

\begin{enumerate}[leftmargin=*, itemsep=2pt]
  \item \textbf{contrôle de la graine}~: un nombre unique, non manquant, fini,
    entier et de valeur absolue au plus \code{.Machine\$integer.max}~; sinon,
    et aussi quand la graine est omise, refus \code{graine\_invalide} du module
    \code{outils}. Il n'y a jamais de graine par défaut~: une graine par défaut
    serait une graine non déclarée~;
  \item \textbf{sauvegarde de l'état de l'appelant}~: \code{.Random.seed} de
    l'environnement global, ou son absence, et le réglage \code{RNGkind()}~;
  \item \textbf{tirage}~: \code{set.seed(graine, kind, normal.kind, sample.kind)}
    sous \code{GENERATEUR}, puis évaluation de \code{expr}~;
  \item \textbf{restauration}, par \code{on.exit()}, donc aussi après une erreur
    dans \code{expr}~: le réglage d'abord, parce que \code{RNGkind()} réécrit
    \code{.Random.seed}, puis \code{.Random.seed} de l'appelant, ou sa
    suppression s'il n'existait pas.
\end{enumerate}

\noindent
Les tirages ne dépendent donc que de la graine~: ni de l'état du générateur
de l'appelant, ni de son réglage, que \code{sous\_graine()} laisse inchangés.
Le simulateur (section~\ref{sec:simulation}) est, à cette date, la seule
fonction du package qui tire un aléa~; il reçoit sa graine en argument et la
consigne, avec \code{GENERATEUR}, dans le jeu simulé.

\medskip
\textbf{Limite.} La réserve de la méthode normale de Box-Muller d'un appelant,
gardée par R hors de \code{.Random.seed}, n'est pas restaurable~; c'est sans
effet sur le moteur, qui tire toujours sous \code{Inversion} (limite reprise
d'\code{outil\_USP}, ADR 0007, point 8.1).

\medskip
\textbf{Aucun aléa hors calcul.} Les fichiers temporaires, employés par
exemple pour calculer une empreinte par \code{tools::md5sum()}, sont créés par
\code{tempfile()}, qui ne touche pas \code{.Random.seed} (constat~: dans une
session neuve, \code{.Random.seed} n'existe ni avant ni après \code{tempfile()}
et \code{tools::md5sum()}~; ADR 0007, point 8.2).

\subsection{Vérifications}
\label{sec:verifications-alea}

Deux fichiers de tests portent ces propriétés (section~\ref{sec:batteries})~:

\begin{itemize}[leftmargin=*, itemsep=2pt]
  \item \code{tests/unitaires/test\_outils.R}~: valeur de \code{GENERATEUR}~;
    mêmes tirages que \code{set.seed()} nu sous le générateur par défaut~;
    graine entière écrite en double admise~; appelant sous
    \code{L'Ecuyer-CMRG}~: tirages de \code{GENERATEUR}, \code{.Random.seed}
    et \code{RNGkind()} de l'appelant restaurés~; absence de \code{.Random.seed}
    restaurée~; état restauré après une erreur dans \code{expr}~; refus
    \code{graine\_invalide} des graines mal formées~;
  \item \code{tests/test\_reproductibilite.R}~: à graine égale, tirages
    identiques (\code{identical()}) dans un même processus, quel que soit l'état
    du générateur de l'appelant, et dans un autre processus R de la même
    machine~; à graines différentes, tirages différents. Les tirages comparés
    couvrent les trois réglages du générateur (\code{runif()}, \code{rnorm()},
    \code{sample.int()})~; le processus fils les écrit en \code{\%.17g}, format
    qui restitue un double au bit près. La même batterie vérifie le simulateur
    et la construction (sections~\ref{sec:simulation} et~\ref{sec:rapport}).
\end{itemize}

\noindent
Ces vérifications sont des constats sur la machine qui exécute la batterie~;
sur la CI, c'est la plateforme de référence (section~\ref{sec:plateforme}).

% =============================================================================
\section{Batteries de vérification et plateforme de CI}
\label{sec:batteries}
% =============================================================================

\subsection{Trois batteries}
\label{sec:trois-batteries}

Les vérifications du code passent par trois batteries en R de base, sans
\code{testthat} (package du CRAN, hors priorité \code{base}), lancées depuis
la racine du dépôt (ADR 0007, point 11.1)~:

\begin{encadre}
\small
\code{Rscript tests/test\_unitaires.R}\\
\code{Rscript tests/test\_reproductibilite.R}\\
\code{Rscript tests/concordance\_doc\_moteur.R -{}-strict}
\end{encadre}

\noindent
La même liste, dans le même ordre, figure dans \code{CLAUDE.md}
(«~Commandes~»), dans le workflow \code{.claude/workflows/circuit-technique.js}
et dans le job «~Tests (R 4.3.3)~» de la CI~; un script le contrôle
(section~\ref{sec:controles-depot}). Chaque batterie~:

\begin{itemize}[leftmargin=*, itemsep=2pt]
  \item \textbf{installe l'arbre de travail} dans une bibliothèque temporaire
    (\code{R CMD INSTALL -{}-no-docs -l <bibliothèque> <racine>}, par
    \code{installer\_tablexp()}), sans installer aucun autre package, puis
    charge \code{tablexp} depuis elle et échoue s'il est chargé d'ailleurs~:
    c'est le code de l'arbre qui est testé, jamais un package installé
    auparavant~;
  \item affiche \code{R.version.string} puis, à la fin, une \textbf{ligne de
    bilan}~: version de R, nombre de tests, d'échecs, d'échecs attendus et de
    succès inattendus~;
  \item \textbf{sort avec un code non nul} en cas d'échec ou de succès
    inattendu (section~\ref{sec:harnais}).
\end{itemize}

\noindent
Une bibliothèque déjà installée, désignée par la variable d'environnement
\code{TABLEXP\_BIBLIOTHEQUE}, n'est réutilisée que si l'\emph{empreinte de
l'arbre source} inscrite à son installation est égale à celle de l'arbre
courant~; sinon la batterie réinstalle et le dit. Cette empreinte est un MD5
(\code{tools::md5sum()}) de \code{DESCRIPTION}, de \code{NAMESPACE} et, s'ils
existent, de tous les fichiers de \code{R/}, \code{inst/}, \code{data/} et
\code{src/}, chemins relatifs compris, triés par octets
(\code{empreinte\_source()}). Elle empêche de tester un package périmé~; elle
n'a aucun rôle dans les résultats.

\medskip
\textbf{\code{tests/test\_unitaires.R}} exécute chaque fichier
\code{tests/unitaires/test\_*.R}, triés par nom, chacun dans son
propre environnement~; \code{quit()}, \code{q()} et \code{bilan()} y sont
neutralisés, une erreur au chargement d'un fichier est un échec, et un fichier
sans aucune assertion aussi. \textbf{\code{tests/test\_reproductibilite.R}}
vérifie les identités de la section~\ref{sec:verifications-alea}, puis celles du
simulateur et de la construction~; la comparaison aux références de
non-régression s'y ajoutera avec la branche qui les crée.
\textbf{\code{tests/concordance\_doc\_moteur.R}} confronte le présent document
au code installé~: toute fonction citée existe~; toute fonction exportée est
citée, ou inscrite dans la liste d'attente \code{tests/concordance\_attente.csv}
avec renvoi à l'issue qui la documentera, où elle compte en échec attendu~;
toute ligne de cette liste désigne un export existant et renvoie à une issue.
Le format des citations et des catalogues qu'elle lit est fixé par
\code{docs/latex/CONVENTIONS.md} (section~3). Sans \code{-{}-strict}, elle
rend compte des écarts sans échouer.

\subsection{Harnais}
\label{sec:harnais}

Les trois batteries partagent \code{tests/harnais.R}, repris de la forme
d'\code{outil\_USP} (ADR 0007, point 11.2, et annotation du 16/10/2026,
point~d)~:

\begin{itemize}[leftmargin=*, itemsep=2pt]
  \item \code{verifier(nom, expr)}~: succès si \code{expr} vaut \code{TRUE}~;
    une chaîne décrit l'écart~; toute autre valeur, ou une erreur, est un
    échec~;
  \item \code{echec\_attendu(nom, issue, expr, motif, erreur\_quelconque)}~:
    \textbf{défaut connu}, codé en échec attendu avec renvoi obligatoire à
    l'issue, de la forme \code{\#N}~; \code{expr} décrit le comportement
    \emph{correct}. Tant qu'elle échoue, le test est «~attendu~»~; \textbf{si
    elle passe, c'est un succès inattendu et la batterie échoue}~: le défaut
    semble corrigé, la marque est retirée et le test devient ordinaire. Trois
    natures d'échec attendu sont admises et décomptées~: par valeur
    (\code{expr} rend \code{FALSE} ou une chaîne), par refus (un
    \code{refus\_tablexp}, du motif donné s'il l'est) et, seulement si
    \code{erreur\_quelconque = TRUE}, par une autre erreur~; toute autre issue
    est un échec, pour qu'un défaut connu ne masque pas une panne différente~;
  \item \code{proche(a, b, rel, abs)}~: égalité numérique à tolérance déclarée,
    $|a-b| \le \mathit{abs} + \mathit{rel}\,|b|$ élément par élément, sans
    tolérance par défaut~; rend un écart pour des valeurs non numériques, des
    longueurs différentes ou des valeurs manquantes à des positions
    différentes~; exige l'égalité exacte des valeurs infinies~;
  \item \code{leve\_refus(expr, motif)}~: vrai si \code{expr} lève un
    \code{refus\_tablexp} du motif donné~;
  \item \code{sans\_execution(x)}~: retire de l'objet résultat le champ des
    informations d'exécution, seul champ qui change entre deux exécutions
    identiques, avant \code{identical()}~;
  \item \code{bilan()}~: ligne de bilan et code de sortie. Il compte en outre
    un échec si un espace de noms hors priorité \code{base}, autre que
    \code{tablexp}, est chargé en fin de batterie (\code{espaces\_hors\_base()})~:
    l'image de la CI contient des packages \code{recommended}, et cette règle
    empêche un test de s'appuyer sur l'un d'eux à l'insu de la liste blanche
    (ADR 0008, «~Conséquences~»).
\end{itemize}

\noindent
Une \textbf{garde de sortie} (\code{.Last()}) fait sortir avec le code~1 tout
processus qui a chargé le harnais et se termine avant la ligne de bilan, même
par \code{quit(status = 0)}. Le harnais est lui-même testé
(\code{tests/unitaires/test\_harnais.R})~: une assertion fausse donne un code
de sortie non nul, et un échec attendu qui réussit fait échouer la batterie.

\medskip
Chaque test nomme la nature de son attendu (ADR 0007, point 11.4~; E2.20)~:
valeur calculée à la main (identité exacte, ou tolérance déclarée)~;
\emph{recoupement}, comparaison à la valeur que rend, dans la batterie, un appel
direct d'une fonction d'un package de priorité \code{base}, dont la portée est
limitée quand le moteur repose sur la même fonction~;
\emph{contre-vérification}, valeur de référence produite hors CI par un package
avancé et écrite en dur avec sa version et sa date~; \emph{constat de
simulation}, qui écrit le nombre de têtes, le niveau, la région d'acceptation
et la probabilité de fausse alarme (\code{CONTEXT.md}~; E2.19).

\subsection{Plateforme de référence}
\label{sec:plateforme}

Les références de non-régression, qui viendront avec la deuxième branche, sont
produites et comparées sur une \emph{plateforme de référence} unique (E4.7~;
ADR 0008)~: le job de la CI s'exécute dans un \textbf{conteneur
\code{rocker/r-ver:4.3.3} épinglé par empreinte}, dont R charge la \textbf{BLAS
et la LAPACK de référence} (Netlib), en monothread. Identifiant fixé par
l'annotation du 16/10/2026 de l'ADR 0008 (décision M47)~:

\begin{encadre}
\small
image \code{rocker/r-ver:4.3.3@sha256:}\\
\code{\phantom{image }732d15020af326da9e919c07f70ca32bf5d3e409220af32e0a4b6d0a89437309}\\
suivie de la bascule \code{update-alternatives} vers les paquets \code{libblas3}
et \code{liblapack3} 3.10.0-2ubuntu1 d'Ubuntu~22.04~;\\
\code{R version 4.3.3 (2024-02-29)}~;\\
BLAS \code{/usr/lib/x86\_64-linux-gnu/blas/libblas.so.3.10.0}~;\\
LAPACK \code{/usr/lib/x86\_64-linux-gnu/lapack/liblapack.so.3.10.0}, version 3.10.0~;\\
monothread (\code{OPENBLAS\_NUM\_THREADS} et \code{OMP\_NUM\_THREADS} valent 1).
\end{encadre}

\noindent
\textbf{Pourquoi un conteneur épinglé.} Un libellé d'exécuteur
(\code{ubuntu-24.04}) n'est pas une image figée~: elle est mise à jour par
GitHub, bibliothèques comprises, sans commit dans le dépôt. \code{outil\_USP},
qui installait R sur l'image de l'exécuteur, a constaté une dérive entre deux
exécutions de même R et de même image, due à une BLAS optimisée (OpenBLAS) qui
choisit ses noyaux selon le processeur (constat d'\code{outil\_USP} rapporté
par l'ADR 0008, «~Contexte~»). Une étiquette
de conteneur peut elle aussi désigner une autre image après une
reconstruction (hypothèse de l'ADR 0008, non vérifiée)~; l'empreinte
(\code{@sha256:\ldots}) désigne un contenu immuable. Le conteneur est une fonctionnalité native de GitHub Actions, et
non une action tierce, ce que la politique des Actions du dépôt exige.

\medskip
\textbf{Pourquoi la bascule.} L'image charge par défaut OpenBLAS (relevé de la
première exécution, rapporté par l'annotation du 16/10/2026 de l'ADR 0008). Le mainteneur a retenu
de basculer dans le job vers la Netlib déjà présente dans l'image, plutôt que
de compiler R depuis la source (décision M47~; le repli par compilation reste
prévu par l'ADR 0008, point 10). L'étape «~BLAS et LAPACK de référence~» de
\code{.github/workflows/ci.yml} pointe par \code{update-alternatives} les
alternatives \code{libblas.so.3-x86\_64-linux-gnu} et
\code{liblapack.so.3-x86\_64-linux-gnu} vers \code{libblas3} et
\code{liblapack3}, puis affiche les cibles résolues et les versions des
paquets. La plateforme de référence est l'image épinglée \emph{suivie de cette
bascule}~: modifier l'étape, l'image ou les bibliothèques est un changement de
plateforme, décidé par le mainteneur et, une fois des références existantes,
accompagné d'un tableau avant / après de toutes les références (ADR 0008,
point 9).

\medskip
\textbf{Script de plateforme.} \code{tests/controles/plateforme.R} s'exécute en
première étape de chaque job R, avant tout calcul (ADR 0008, points 3 à 5).
En mode strict (\code{-{}-exiger}, sur la CI), il \textbf{échoue} si~:
\code{R.version.string} n'est pas exactement \code{R\_ATTENDU}
(\code{R version 4.3.3 (2024-02-29)})~; l'une des quatre lectures de la BLAS et
de la LAPACK (\code{extSoftVersion()["BLAS"]}, \code{sessionInfo()\$BLAS},
\code{sessionInfo()\$LAPACK}, \code{La\_library()}) est vide, désigne une BLAS
optimisée (motif \code{MOTIF\_OPTIMISE}~: \code{openblas}, \code{mkl},
\code{atlas}, \code{blis}, \code{flexiblas}) ou sort des chemins de référence
(\code{MOTIF\_BLAS}, \code{MOTIF\_LAPACK})~; \code{La\_version()} n'est pas
\code{LA\_VERSION\_ATTENDUE} (3.10.0)~; \code{OPENBLAS\_NUM\_THREADS} ou
\code{OMP\_NUM\_THREADS} ne vaut pas 1~; \code{TABLEXP\_IMAGE}, que le job
renseigne avec la même chaîne que l'image du conteneur, n'est pas épinglée par
empreinte (\code{MOTIF\_IMAGE}). Sans \code{-{}-exiger} (poste, session), il
affiche les mêmes grandeurs et les écarts sans échouer~: les batteries restent
exécutables sur le poste Windows du mainteneur, dont la BLAS est celle de R.
Dans les deux modes, il écrit l'\emph{identifiant de plateforme} (image,
\code{R.version.string}, BLAS, LAPACK et sa version, \code{sessionInfo()}) et
consigne sans faire échouer le processeur et ses drapeaux (\code{fma},
\code{avx2}, \code{avx512f}), \code{ImageOS} et \code{ImageVersion} de l'hôte,
\code{LC\_COLLATE}, le fuseau horaire, la présence d'un exécutable
\code{java} et \code{GITHUB\_SHA}. Ce bloc sera reporté dans le message de
chaque commit de référence.

\begin{encadre}
\small
\textbf{Ce que la plateforme ne garantit pas.} Le modèle de processeur varie
d'une exécution à l'autre sous le même libellé d'exécuteur (relevé de l'ADR
0008, annotation du 16/10/2026). La BLAS Netlib n'a pas de noyau choisi selon
le processeur~; la réserve sur la bibliothèque mathématique du système
(section~\ref{sec:reproductibilite-exigence}) reste en revanche ouverte, non
mesurée. Les références seront donc comparées à tolérance, structure contrôlée
à part (E4.7).
\end{encadre}

\subsection{Jobs de la CI}
\label{sec:jobs-ci}

\code{.github/workflows/ci.yml} définit les jobs suivants, exécutés sur chaque
pull request, sur chaque push vers \code{main} et sur demande. Aucun job n'installe de
package du CRAN ni ne configure de dépôt CRAN pour une installation~:
\code{tablexp} est installé depuis l'arbre de travail par chaque batterie.

\begin{center}
\small
\rowcolors{2}{tablerow}{white}
\begin{longtable}{@{}p{3.3cm}p{3.2cm}p{8.3cm}@{}}
\caption{Jobs de la CI (\code{.github/workflows/ci.yml}).}
\label{tab:jobs-ci}\\
\rowcolor{tableheader}
\textcolor{financeblue}{\textbf{Job}} &
\textcolor{financeblue}{\textbf{Exécution}} &
\textcolor{financeblue}{\textbf{Contenu}} \\
\endfirsthead
\rowcolor{tableheader}
\textcolor{financeblue}{\textbf{Job}} &
\textcolor{financeblue}{\textbf{Exécution}} &
\textcolor{financeblue}{\textbf{Contenu}} \\
\endhead
«~Contrôles du dépôt~» & hôte \code{ubuntu-24.04}, sans R & réglages de Claude
Code valides (JSON)~; fiches d'agents complètes~; fiches \code{-approfondi} à
jour~; aucun fichier sensible versionné \\
«~Tests (R 4.3.3)~» & conteneur de la plateforme de référence & bascule vers
la BLAS et la LAPACK de référence~; script de plateforme en mode strict~; les
trois batteries~; liste blanche de \code{DESCRIPTION}~; ordre des modules~;
identité des batteries et de l'image. Les batteries et les contrôles
s'exécutent même si une étape antérieure échoue (condition
\code{\$\{\{ !cancelled() \}\}}), pour que tous les écarts soient visibles~;
le job reste en échec \\
«~Absence de Java~» & même conteneur, réseau requis & contrôle de l'absence de
Java (section~\ref{sec:absence-java}) et son cas témoin \\
\end{longtable}
\end{center}

\noindent
Ces trois jobs sont les contrôles requis que liste le ruleset de \code{main}
(\code{.github/ruleset-main.json}, appliqué au dépôt par le mainteneur)~: une
pull request ne peut alors être fusionnée que s'ils sont verts. L'ADR 0008 (point 6) prévoit un quatrième job,
«~Documentation~», qui compilera le PDF du présent document sans action tierce
et le publiera en artefact~; il n'existe pas encore. Les batteries et les
scripts des jobs conteneurisés ne dépendent pas de git, absent de l'image~;
un contrôle qui lit l'historique reste dans «~Contrôles du dépôt~», sur l'hôte
(ADR 0008, point 11).

% =============================================================================
\section{Contrôles du dépôt}
\label{sec:controles-depot}
% =============================================================================

\subsection{Scripts de contrôle}
\label{sec:scripts-controle}

Tout contrôle mécanique est un script, jamais une relecture (\code{CLAUDE.md},
«~Workflows~», principe~7). Les scripts sont en R de base, sous
\code{tests/controles/}, et sortent avec un code non nul en cas d'écart. Chacun
est testé sur des témoins dans \code{tests/unitaires/test\_controles.R}~: il
doit échouer sur un cas fautif et passer sur le dépôt.

\begin{center}
\small
\rowcolors{2}{tablerow}{white}
\begin{longtable}{@{}p{4.6cm}p{7.4cm}p{2.8cm}@{}}
\caption{Scripts de contrôle (\code{tests/controles/}).}
\label{tab:scripts-controle}\\
\rowcolor{tableheader}
\textcolor{financeblue}{\textbf{Script}} &
\textcolor{financeblue}{\textbf{Règle contrôlée}} &
\textcolor{financeblue}{\textbf{Source}} \\
\endfirsthead
\rowcolor{tableheader}
\textcolor{financeblue}{\textbf{Script}} &
\textcolor{financeblue}{\textbf{Règle contrôlée}} &
\textcolor{financeblue}{\textbf{Source}} \\
\endhead
\code{liste\_blanche\_}\newline\code{description.R} & \code{Depends}~: exactement
\code{R (>= 4.3.3)}, plus d'éventuels packages de priorité \code{base} ou
admis~; \code{Imports}~: packages de priorité \code{base} de R 4.3.3, ou
admis~; \code{LinkingTo}, \code{Suggests}, \code{Enhances}~: packages admis
seulement. Un package est admis s'il figure, avec son champ et l'ADR de
méthode qui l'introduit, dans la constante \code{ADMIS}, vide à cette date &
ADR 0002, point 3 et A8~; M44 \\
\code{graphe\_modules.R} & graphe d'appels autorisés et appels dynamiques
interdits (section~\ref{sec:graphe}) & ADR 0007, points 2.1 à 2.3 \\
\code{identite\_batteries.R} & mêmes batteries, dans le même ordre, dans
\code{CLAUDE.md}, dans \code{circuit-technique.js} et dans les étapes
«~Batterie~:~» du job «~Tests (R 4.3.3)~»~; ce job ne porte ni
\code{continue-on-error}, ni condition \code{if:} de job, et une étape de
batterie n'a pas d'autre condition que \code{\$\{\{ !cancelled() \}\}}~; une
seule chaîne d'image dans tous les workflows & \code{CLAUDE.md}~; ADR 0008,
point 1 \\
\code{plateforme.R} & plateforme de référence (section~\ref{sec:plateforme}) &
ADR 0008, points 3 à 5 \\
\code{absence\_java.R} & aucune dépendance à Java
(section~\ref{sec:absence-java}) & ADR 0002, A1~; ADR 0008, point 7~; issue~\#3 \\
\end{longtable}
\end{center}

\noindent
La liste des packages de priorité \code{base} est écrite en dur dans
\code{liste\_blanche\_description.R} pour R 4.3.3 (\code{PACKAGES\_BASE\_R433})~:
le verdict ne dépend pas de la machine~; le script signale, sans échouer, un
écart avec les packages de priorité \code{base} effectivement installés.

\subsection{Absence de Java}
\label{sec:absence-java}

\textbf{Règle.} Aucun code du projet ne dépend de Java, ni directement ni
indirectement~: ni le moteur, ni l'application, ni la contre-vérification
(ADR 0002, annotation du 16/10/2026, A1, décision M40). Java n'est pas
utilisable sur le poste de travail du mainteneur~; c'est pour cette raison
que le package \code{ELT}, qui dépend de \code{xlsx} et donc de Java, est exclu
de la contre-vérification. Sa source ne peut être exécutée qu'hors du dépôt,
sans Java, en appoint, sans qu'aucun code en soit copié (ADR 0002, A5 et
troisième annotation).

\medskip
\textbf{Contrôle.} \code{tests/controles/absence\_java.R} (issue~\#3) contrôle
trois ensembles de packages~:

\begin{itemize}[leftmargin=*, itemsep=2pt]
  \item le \emph{moteur}~: les packages déclarés dans les champs \code{Depends},
    \code{Imports} et \code{LinkingTo} de \code{DESCRIPTION}, hors R~; aucun à
    cette date~;
  \item l'\emph{application}~: \code{shiny}, \code{htmltools}, et
    \code{openxlsx} à titre facultatif~;
  \item la \emph{contre-vérification}~: les packages \code{survival},
    \code{mgcv}, \code{WH}, \code{StMoMo}, \code{demography}, \code{MortalityLaws},
    \code{lifecontingencies}, \code{survPen}, \code{MortCast}, \code{locfit}, et
    \code{flexsurv} à titre facultatif.
\end{itemize}

\noindent
Les deux derniers ensembles sont lus dans une liste unique, le fichier
\code{packages\_hors\_moteur.csv} du dossier \code{tests/controles/} (colonnes \code{ensemble},
\code{package}, \code{statut}, \code{source}), qui est la référence~; la liste
ci-dessus en est la transcription à la date du document. Le script lit les
métadonnées d'un \textbf{instantané daté} du CRAN chez Posit Package Manager,
par \code{available.packages()}, sans rien installer~; la date est la constante
\code{INSTANTANE} du script (2026-10-15, un jour ouvré~: selon la documentation de Posit citée par
l'ADR 0008, aucun instantané n'est publié les week-ends et jours fériés), qui ne change que par un commit
\code{repo:}. Il calcule, pour chaque package, la fermeture récursive de ses
dépendances \code{Depends}, \code{Imports} et \code{LinkingTo}, par la
fonction \code{tools::package\_dependencies()} appelée avec
\code{recursive = TRUE}, et \textbf{échoue} (code~1)~:

\begin{itemize}[leftmargin=*, itemsep=2pt]
  \item si un package de \code{PACKAGES\_JAVA} (\code{rJava}, \code{xlsx},
    \code{xlsxjars}) figure dans un ensemble ou dans la fermeture de l'un de ses
    packages~; la chaîne de dépendances la plus courte qui y mène est
    affichée~;
  \item si un package listé, ou une dépendance récursive hors priorité
    \code{base}, manque aux métadonnées~: un package introuvable n'a pas de
    dépendances connues, il ne passe pas en silence~;
  \item si la liste est mal formée (en-tête, ensembles et statuts connus, noms
    valides, sans doublon, ensembles non vides).
\end{itemize}

\noindent
Des métadonnées illisibles (panne réseau, adresse fausse) donnent le code~2~:
aucun verdict n'est rendu, l'échec du job n'est jamais masqué et une relance
le lève. Seul le filtre \code{duplicates} d'\code{available.packages()} est
appliqué~: le filtre par version de R écarterait les packages qui exigent un R
plus récent que celui qui exécute le contrôle, alors que la
contre-vérification peut tourner sur un autre poste muni d'un R récent, et ces
packages seraient alors signalés absents (faux échec)~; le filtre par système
d'exploitation est écarté aussi, le poste étant sous Windows et la CI sous
Linux.

\medskip
\textbf{Cas témoin.} Le job «~Absence de Java~» exécute aussi le script avec
\code{-{}-ajouter ELT}~: il doit rendre le code~1 (un verdict, et non le code~2
d'une panne) et afficher la chaîne de \code{ELT} à \code{xlsx}. Les tests du
script s'exécutent sans réseau, sur des dépôts locaux témoins
(\code{tests/unitaires/test\_controles.R}).

\begin{encadre}
\small
\textbf{Limite.} La détection porte sur les seuls noms de
\code{PACKAGES\_JAVA}. Le champ \code{SystemRequirements}, qui pourrait déclarer
Java, n'est pas renseigné dans les métadonnées \code{PACKAGES} du CRAN (relevé
consigné dans l'en-tête du script)~: un package qui exigerait Java sans
dépendre de l'un de ces noms échapperait au contrôle. Étendre la détection
relève du mainteneur.
\end{encadre}


% =============================================================================
\partie{Partie II}{Point d'entrée et plan de construction}
% =============================================================================

\section{Point d'entrée unique~: \code{R/construire\_table.R}}
\label{sec:construction}

\subsection{Plan de construction et validation}
\label{sec:plan}

\emph{À rédiger (issue~\#11).}

\subsection{Objet résultat}
\label{sec:objet-resultat}

\emph{À rédiger (issue~\#11).}

\subsection{Fonctions du module}
\label{sec:fonctions-construire-table}

\emph{À rédiger (issue~\#11).}


% =============================================================================
\partie{Partie III}{Module \code{donnees}}
% =============================================================================

\section{Données, rôles et contrôles~: \code{R/donnees.R}}
\label{sec:donnees}

\subsection{Catalogue des rôles}
\label{sec:catalogue-roles}

\emph{À rédiger (issue~\#9).}

\subsection{Contrôles, refus et exclusions}
\label{sec:controles-donnees}

\emph{À rédiger (issue~\#9).}

\subsection{Unité d'observation et fusion par tête}
\label{sec:unite-observation}

\emph{À rédiger (issue~\#9).}

\subsection{Table de référence}
\label{sec:table-reference}

\emph{À rédiger (issue~\#9).}

\subsection{Fonctions du module}
\label{sec:fonctions-donnees}

\emph{À rédiger (issue~\#9).}


% =============================================================================
\partie{Partie IV}{Module \code{exposition}}
% =============================================================================

\section{Exposition et événements~: \code{R/exposition.R}}
\label{sec:exposition}

\subsection{Recette d'exposition et conventions d'âge}
\label{sec:recettes}

\emph{À rédiger (issue~\#11).}

\subsection{Troncature à gauche et censure à droite}
\label{sec:troncature-censure}

\emph{À rédiger (issue~\#11).}

\subsection{Découpage par année}
\label{sec:decoupage-annee}

\emph{À rédiger (issue~\#11).}

\subsection{Fonctions du module}
\label{sec:fonctions-exposition}

\emph{À rédiger (issue~\#11).}


% =============================================================================
\partie{Partie V}{Module \code{brut}}
% =============================================================================

\section{Taux bruts~: \code{R/brut.R}}
\label{sec:brut}

Module prévu (E2.12, E2.14)~: estimateurs des taux bruts et leurs intervalles
de confiance exacts et asymptotiques. Il sera créé par la branche des taux
bruts (\code{docs/}\allowbreak\code{feuille-de-route.md}, §~2).

\subsection{Fonctions du module}
\label{sec:fonctions-brut}

{\small
\rowcolors{2}{tablerow}{white}
\begin{longtable}{@{}p{4.2cm}p{1.6cm}p{3.4cm}p{5.4cm}@{}}
\caption{Définitions de premier niveau de \code{R/brut.R}.}
\label{tab:fonctions-brut}\\
\rowcolor{tableheader}
\textcolor{financeblue}{\textbf{Nom}} &
\textcolor{financeblue}{\textbf{Statut}} &
\textcolor{financeblue}{\textbf{Entrées}} &
\textcolor{financeblue}{\textbf{Sortie et rôle}} \\
\multicolumn{4}{@{}l@{}}{aucune fonction~: fichier \code{R/brut.R} absent} \\
\end{longtable}
}


% =============================================================================
\partie{Partie VI}{Module \code{lissage}}
% =============================================================================

\section{Lissage~: \code{R/lissage.R}}
\label{sec:lissage}

Module prévu (E2.12)~: lissages des taux bruts. Il sera créé par une branche
ultérieure (\code{docs/}\allowbreak\code{feuille-de-route.md}, §~2).

\subsection{Fonctions du module}
\label{sec:fonctions-lissage}

{\small
\rowcolors{2}{tablerow}{white}
\begin{longtable}{@{}p{4.2cm}p{1.6cm}p{3.4cm}p{5.4cm}@{}}
\caption{Définitions de premier niveau de \code{R/lissage.R}.}
\label{tab:fonctions-lissage}\\
\rowcolor{tableheader}
\textcolor{financeblue}{\textbf{Nom}} &
\textcolor{financeblue}{\textbf{Statut}} &
\textcolor{financeblue}{\textbf{Entrées}} &
\textcolor{financeblue}{\textbf{Sortie et rôle}} \\
\multicolumn{4}{@{}l@{}}{aucune fonction~: fichier \code{R/lissage.R} absent} \\
\end{longtable}
}


% =============================================================================
\partie{Partie VII}{Module \code{positionnement}}
% =============================================================================

\section{Positionnement sur une table de référence~: \code{R/positionnement.R}}
\label{sec:positionnement}

Module prévu (E2.12, E2.15)~: positionnement d'une table d'expérience sur une
table de référence. Il sera créé par une branche ultérieure
(\code{docs/}\allowbreak\code{feuille-de-route.md}, §~2).

\subsection{Fonctions du module}
\label{sec:fonctions-positionnement}

{\small
\rowcolors{2}{tablerow}{white}
\begin{longtable}{@{}p{4.2cm}p{1.6cm}p{3.4cm}p{5.4cm}@{}}
\caption{Définitions de premier niveau de \code{R/positionnement.R}.}
\label{tab:fonctions-positionnement}\\
\rowcolor{tableheader}
\textcolor{financeblue}{\textbf{Nom}} &
\textcolor{financeblue}{\textbf{Statut}} &
\textcolor{financeblue}{\textbf{Entrées}} &
\textcolor{financeblue}{\textbf{Sortie et rôle}} \\
\multicolumn{4}{@{}l@{}}{aucune fonction~: fichier \code{R/positionnement.R} absent} \\
\end{longtable}
}


% =============================================================================
\partie{Partie VIII}{Module \code{prospective}}
% =============================================================================

\section{Tables prospectives~: \code{R/prospective.R}}
\label{sec:prospective}

Module prévu (E2.13, v2)~: tables prospectives, toute table prospective étant
une surface de mortalité.

\subsection{Fonctions du module}
\label{sec:fonctions-prospective}

{\small
\rowcolors{2}{tablerow}{white}
\begin{longtable}{@{}p{4.2cm}p{1.6cm}p{3.4cm}p{5.4cm}@{}}
\caption{Définitions de premier niveau de \code{R/prospective.R}.}
\label{tab:fonctions-prospective}\\
\rowcolor{tableheader}
\textcolor{financeblue}{\textbf{Nom}} &
\textcolor{financeblue}{\textbf{Statut}} &
\textcolor{financeblue}{\textbf{Entrées}} &
\textcolor{financeblue}{\textbf{Sortie et rôle}} \\
\multicolumn{4}{@{}l@{}}{aucune fonction~: fichier \code{R/prospective.R} absent} \\
\end{longtable}
}


% =============================================================================
\partie{Partie IX}{Module \code{fermeture}}
% =============================================================================

\section{Fermeture aux âges extrêmes~: \code{R/fermeture.R}}
\label{sec:fermeture}

Module prévu (E2.12)~: fermeture de la table aux âges extrêmes. Il sera créé
par une branche ultérieure (\code{docs/}\allowbreak\code{feuille-de-route.md}, §~2).

\subsection{Fonctions du module}
\label{sec:fonctions-fermeture}

{\small
\rowcolors{2}{tablerow}{white}
\begin{longtable}{@{}p{4.2cm}p{1.6cm}p{3.4cm}p{5.4cm}@{}}
\caption{Définitions de premier niveau de \code{R/fermeture.R}.}
\label{tab:fonctions-fermeture}\\
\rowcolor{tableheader}
\textcolor{financeblue}{\textbf{Nom}} &
\textcolor{financeblue}{\textbf{Statut}} &
\textcolor{financeblue}{\textbf{Entrées}} &
\textcolor{financeblue}{\textbf{Sortie et rôle}} \\
\multicolumn{4}{@{}l@{}}{aucune fonction~: fichier \code{R/fermeture.R} absent} \\
\end{longtable}
}


% =============================================================================
\partie{Partie X}{Module \code{validation}}
% =============================================================================

\section{Validation et suivi d'expérience~: \code{R/validation.R}}
\label{sec:validation}

Module prévu (E2.12, E2.18)~: validation de la table construite et suivi
d'expérience. Il sera créé par une branche ultérieure
(\code{docs/}\allowbreak\code{feuille-de-route.md}, §~2).

\subsection{Fonctions du module}
\label{sec:fonctions-validation}

{\small
\rowcolors{2}{tablerow}{white}
\begin{longtable}{@{}p{4.2cm}p{1.6cm}p{3.4cm}p{5.4cm}@{}}
\caption{Définitions de premier niveau de \code{R/validation.R}.}
\label{tab:fonctions-validation}\\
\rowcolor{tableheader}
\textcolor{financeblue}{\textbf{Nom}} &
\textcolor{financeblue}{\textbf{Statut}} &
\textcolor{financeblue}{\textbf{Entrées}} &
\textcolor{financeblue}{\textbf{Sortie et rôle}} \\
\multicolumn{4}{@{}l@{}}{aucune fonction~: fichier \code{R/validation.R} absent} \\
\end{longtable}
}


% =============================================================================
\partie{Partie XI}{Module \code{rapport}}
% =============================================================================

\section{Fiche de construction et empreintes~: \code{R/rapport.R}}
\label{sec:rapport}

\subsection{Fiche de construction}
\label{sec:fiche}

\emph{À rédiger (issue~\#11).}

\subsection{Empreinte des données}
\label{sec:empreinte}

\emph{À rédiger (issue~\#11).}

\subsection{Fonctions du module}
\label{sec:fonctions-rapport}

\emph{À rédiger (issue~\#11).}


% =============================================================================
\partie{Partie XII}{Module \code{simulation}}
% =============================================================================

\section{Simulateur~: \code{R/simulation.R}}
\label{sec:simulation}

\subsection{Loi vraie et tables synthétiques}
\label{sec:loi-vraie}

\emph{À rédiger (issue~\#10).}

\subsection{Portefeuille simulé}
\label{sec:portefeuille-simule}

\emph{À rédiger (issue~\#10).}

\subsection{Constats de simulation}
\label{sec:constats-simulation}

\emph{À rédiger (issue~\#10).}

\subsection{Fonctions du module}
\label{sec:fonctions-simulation}

\emph{À rédiger (issue~\#10).}


% =============================================================================
\partie{Partie XIII}{Index et traçabilité}
% =============================================================================

\section{Index et traçabilité}
\label{sec:index}

\subsection{Index des fonctions exportées}
\label{sec:index-exports}

\emph{À rédiger (issues~\#9, \#10 et \#11).}

\subsection{Traçabilité des exigences}
\label{sec:tracabilite}

Le tableau~\ref{tab:tracabilite} relie les exigences du cahier des charges
(\code{docs/exigences.md}) aux sections qui les décrivent. Il est complété à
chaque section rédigée.

{\small
\rowcolors{2}{tablerow}{white}
\begin{longtable}{@{}p{1.6cm}p{8.6cm}p{4.6cm}@{}}
\caption{Exigences et sections qui les décrivent.}
\label{tab:tracabilite}\\
\rowcolor{tableheader}
\textcolor{financeblue}{\textbf{Exigence}} &
\textcolor{financeblue}{\textbf{Objet}} &
\textcolor{financeblue}{\textbf{Section}} \\
\endfirsthead
\rowcolor{tableheader}
\textcolor{financeblue}{\textbf{Exigence}} &
\textcolor{financeblue}{\textbf{Objet}} &
\textcolor{financeblue}{\textbf{Section}} \\
\endhead
E1.3 & rigueur, traçabilité, reproductibilité &
\ref{sec:reproductibilite}, \ref{sec:batteries} \\
E3.1 & document LaTeX unique, une partie par module, charte d'\code{outil\_USP},
tableaux de fonctions & ensemble du document~; \ref{sec:fonctions-outils} \\
E3.3 & test de concordance entre la documentation et le code &
\ref{sec:trois-batteries} \\
E3.4 & aucun chiffre issu d'une exécution & Objet du document \\
E4.1 & R 4.3.3, package \code{tablexp}, modules & \ref{sec:fichiers-package} \\
E4.2 & packages de priorité \code{base}~; aucune dépendance à Java &
\ref{sec:scripts-controle}, \ref{sec:absence-java} \\
E4.3 & frontière calcul / affichage & \ref{sec:separation} \\
E4.4 & le moteur reçoit un \code{data.frame} et des rôles & \ref{sec:separation} \\
E4.6 & reproductibilité, graine locale, générateur fixé &
\ref{sec:reproductibilite} \\
E4.7 & références de non-régression sur la plateforme de référence &
\ref{sec:plateforme} \\
E4.9 & un module par famille, index des fonctions contrôlé par la
concordance & \ref{sec:architecture}, \ref{sec:trois-batteries} \\
\end{longtable}
}

\subsection{Références}
\label{sec:references}

\begin{itemize}[leftmargin=*, itemsep=3pt]
  \item M.~Matsumoto, T.~Nishimura, \emph{Mersenne Twister: a
    623-dimensionally equidistributed uniform pseudo-random number generator},
    \emph{ACM Transactions on Modeling and Computer Simulation}, 8, 1998,
    p.~3--30 (cité par la documentation de R, \code{?Random})~;
    section~\ref{sec:generateur}.
  \item R Core Team, \emph{R: A Language and Environment for Statistical
    Computing}, R Foundation for Statistical Computing, Vienne, 2024
    (\url{https://www.R-project.org/}~; référence rendue par
    \code{citation()} sous R 4.3.3).
\end{itemize}

\end{document}
